\begin{titlepage}
\hspace{-1.7cm} %Este comando es para mandar a la izquierda :)
%Aquí empiezan los tres minipages de antes
\begin{minipage}[t][0.03\textheight][c]{0.22\textwidth}
        \includegraphics[width=3.7cm]{logo-unach.png}
\end{minipage}\hspace{0.9cm}
\begin{minipage}[t][0.03\textheight][c]{0.69\textwidth}
\begin{center}
                \textsc{\huge Universidad Autónoma de Chiapas}\\[0.3cm]
                \hrule height 2.5pt
                \vspace{0.2cm}
                \hrule height1pt
                \vspace{0.3cm}
                \textsc{\Large Facultad de Ciencias en Física y Matemáticas}
\end{center}
\end{minipage}\hspace{0.2cm}
\begin{minipage}[t][0.03\textheight][c]{0.2\textwidth}
		\includegraphics[width=2.7cm]{FCFMLOGO.png}
\end{minipage}\\
%%%%%%%%%%%%%%%%%%Aquí comienzan los otros dos minipages del título%%%%%%%%%%%%%%%%%%%%%%%%%%%%%%%%%%%%%%%%%%%%%%%%%%%%%%%%%%%%%%%%%%%%%%%%%%%%%%%%%%%%%%%%%%%%%%
\begin{minipage}[t][0.93\textheight][c]{0.06\textwidth}
\vspace{60pt}
    \begin{center}
        \vrule width1pt height18cm
        \vspace{5mm}
        \vrule width2.5pt height18cm
        \vspace{5mm}
        \vrule width1pt height18cm
   \end{center}
\end{minipage}\hspace{1.3cm} %Esto mueve las letras que tienes ahí
\begin{minipage}[t][0.95\textheight][c]{0.76\textwidth}

            \begin{center}
                \vspace{0.7cm}
                {\Large\bfseries AQUI VA EL NOMBRE DE LA TESIS}\\[2cm]
                \textsc{\huge \textbf{T\, E\, S\, I\, S}}\\[1.5cm]
                \textsc{\large QUE PARA OBTENER EL TÍTULO DE:}\\[0.3cm]
                \textbf{\textsc{LICENCIADO EN MATEMÁTICAS APLICADAS}}\\[1.5cm]
                \textsc{\large PRESENTA:}\\[0.3cm]
                \textbf{\textsc{\large {CRISTÓBAL PÉREZ GONZÁLEZ}}}\\[2cm]
                {\large\scshape Director:\\[0.3cm]
                {\textbf{\large Dr. Yofre Hernán García Gómez }}}\\[2.0cm]
                \large{Tuxtla Gutiérrez, Chiapas a \today.}

            \end{center}
\end{minipage}
\end{titlepage}

\pagebreak[2]

\chapter*{Dedicatoria}
\begin{flushright}
\itshape
\setlength{\parindent}{0pt} % Elimina sangría para mayor elegancia
\setlength{\parskip}{1.2\baselineskip} % Espacio generoso entre párrafos

A Dios, por guiar mis pasos y darme la fortaleza para transformar la distancia en propósito.

A mis padres, Nelva y Alfonso, por sostenerme desde lejos con llamadas al amanecer, por su apoyo incondicional, sus sacrificios silenciosos y por creer en mí incluso cuando yo misma dudaba.

A mi hermanita Esmeralda, por recordarme con su sonrisa que existía un mundo más allá de las ecuaciones y por ser mi puente a casa en cada audio y cada foto.

Y a toda mi familia Vázquez, por convertir cada visita en refugio y cada despedida en promesa.

Este título no solo lleva mi nombre: lleva el peso amoroso de quienes me esperaron al otro lado del camino.
\end{flushright}


\chapter*{Agradecimientos}

En primer lugar, quiero agradecer a Dios por toda la sabiduría que me ha brindado, tanto durante la escritura de esta tesis y mi paso por la carrera, como a lo largo de toda mi vida. Gracias por darme las fuerzas necesarias cuando más las necesitaba y en aquellos momentos en los que sentía que las cosas no me iban a salir.

Quiero expresar mi más sincero agradecimiento al Dr. Yofre Hernán García Gómez, mi director de tesis, por todo el apoyo y la confianza que ha tenido en mí. Aunque esta investigación tuvo varios momentos difíciles en los que creíamos que algo no nos saldría, gracias a su guía y a su gran paciencia conmigo nunca me di por vencido. Finalmente, le agradezco por enseñarme que las Matemáticas Aplicadas tienen la capacidad de ayudar y mejorar la vida de las personas.

Asimismo, agradezco a todos mis profesores y profesoras que tuve a lo largo de la licenciatura. Al Dr. Aarón, al Dr. Sergio, al Dr. Javier y al Dr. Saúl, gracias por enseñarme de la mejor forma y con el mayor rigor toda la matemática que hoy sé. Al Dr. Boris, por ser mi tutor durante toda la carrera; gracias por creer en mí desde el primer momento y por animarme siempre a seguir adelante en el hermoso camino de las Ciencias Exactas. Al Mtro. Elesban, por su forma tan especial de enseñar y por hacer que amara la programación. A la Mtra. Marda, por ser más que una maestra, una amiga; gracias por sus regaños que me ayudaron a mejorar como persona y por ser quien me introdujo al mundo de las ecuaciones diferenciales. De igual manera, quiero agradecer al Dr. Yofre y al Dr. Armando por su forma tan pacífica de enseñar y por haberme introducido al hermoso mundo de la probabilidad y de los procesos estocásticos.

De manera muy especial, agradezco a mis padres, María Teresa y Sebastián, por todo el apoyo que me han brindado a lo largo de mi vida, tanto en lo económico como en lo emocional. Gracias por impulsarme siempre a seguir adelante, por nunca dudar de mí y por motivarme en todo momento a estudiar esta carrera. También agradezco a mis hermanos, Saraí, Cristian y Cristóbal, por estar conmigo, motivarme y hacerme compañía cada vez que estoy en casa con ellos.

Por último, pero no menos importante, quiero agradecer a mis amigos. A Carlos y Marbin, por ser las personas con las que más conviví a lo largo de la carrera, por acompañarme en este proceso tan largo pero tan gratificante, y por estar siempre conmigo cuando más los necesitaba o cuando me sentía sin ánimos. A Mariana, Michelle y Sofía, por estar a mi lado en mis momentos más difíciles; gracias por todos sus consejos, sus ánimos y los abrazos que me dieron cuando sentía que ya no podía más. Gracias a todos ustedes por ser los mejores amigos que pude tener en mi vida.

A todos ustedes, muchas gracias, ya que sin su apoyo incondicional, sus consejos y su motivación, este logro no lo habría podido conseguir yo solo.